%! Author = chtompki
%! Date = 1/14/26
\documentclass[12pt]{amsart}
% Default margins are too wide all the way around. I reset them here
\setlength{\hoffset}{-.75in}
\setlength{\textwidth}{6.5in}
\setlength{\voffset}{-.5in}
\setlength{\textheight}{9.0in}
\usepackage{amsmath}
\usepackage{amssymb}
\usepackage{amsthm}
\usepackage{enumitem}
\usepackage{setspace}
\usepackage{graphicx}

\newenvironment{solution}
{\begin{proof}[Solution]}
{\end{proof}}

\NewDocumentCommand{\vitem}{v}{%
    \item[\texttt{#1}]%
}

\title{Numerical Analysis\\Problem Set 2}
\author{Rob Tompkins}
\date{\today}

\begin{document}
    \maketitle
    \date{\today}
    \begin{onehalfspace}
        We begin by noting that all of the files referenced in this document are available at: \\
        \underline{https://github.com/chtompki/MathSpring2026/tree/main/NumericalAnalysis/ProblemSet2}
        % Problem 1
        \textbf{1.} Use the command \texttt{magic(n)}, for $n=100,200,300,\ldots,10^4$ to construct the magic square
        matrices of order $n$.
        Magic squares are defined as those matrices having entries for which the sum of the elements by rows, columns,
        or diagonals are identical.
        Then compute their determinants by the command \texttt{det()}.
        Use the command \texttt{cputime} to get the CPU time that is needed fore this computation.
        \bf{Be patient, it may take around 10 minutes to compute these determinants.}
    \end{onehalfspace}
\end{document}